\chapter{The Supervisor Portal}
\label{ch:supervisor}

\portalcover{Supervisor Portal}{For the person in charge of a department}{secSort}

\section{Who this chapter is for}

You run a department. You can do everything an operator can do at a station ---
receiving, dispatching, scanning, tagging --- and four things they cannot:

\begin{bullets}
  \item See how much work your section is holding and how long it has been there.
  \item \textbf{Settle a count difference} that the operators cannot resolve.
  \item \textbf{Cancel a handover} that was sent by mistake.
  \item Build and save reports for your own use and for the team.
\end{bullets}

\section{Checking the state of your section}

Open \menu{Section WIP} and choose your department at the top.

\shot{12-admin-sections-wip}{Section WIP. The four boxes summarise the section;
the table below breaks the work down however you choose.}{fig:sup-wip}

\begin{table}[H]
\centering
\small
\begin{tabular}{@{}p{32mm}p{106mm}@{}}
\toprule
\textbf{Box} & \textbf{What it tells you} \\
\midrule
Garments in process & How much work is physically in your section right now. \\
Groups              & How many different kinds of work that is. \\
Average wait        & How long a typical piece has been sitting. \\
Longest wait        & The oldest piece. This is the number to watch. \\
\bottomrule
\end{tabular}
\end{table}

\subsection{Slicing the work however you need}

The \emph{Group by} row lets you break the same work down in different ways. Tick
any combination:

\begin{table}[H]
\centering
\small
\begin{tabular}{@{}p{34mm}p{104mm}@{}}
\toprule
\textbf{Group by} & \textbf{Answers the question} \\
\midrule
Design / Colour / Size & What styles am I holding? \\
Customer / Order       & Whose goods am I holding? \\
Receiving batch        & \textbf{Which delivery did this come in on, and when?} \\
Age bucket             & How much of my work is fresh and how much is stale? \\
Article status         & Is anything on hold or in rework? \\
Fabric type            & Useful when a shade problem is suspected. \\
\bottomrule
\end{tabular}
\end{table}

\shot{13-admin-sections-batch}{The same work grouped by the delivery it arrived
on. The panel underneath lists every bulk receipt into the section, with the time
it was received and how much of it is still here.}{fig:sup-batch}

\begin{tipbox}
\textbf{Finding stalled work.} Set \field{Waiting at least (hours)} to 24 and set
\field{Sort by} to ``Oldest arrival first''. What you are left with is exactly
the work that is holding up your section.
\end{tipbox}

Press \btn{Export CSV} to take any view into Excel.

\section{Settling a count difference}

When operators have scanned a batch twice and pieces are still missing, only a
supervisor can close it.

\begin{steps}
  \item Open \menu{Transfers} and find the document. Its status is
        \texttt{VARIANCE}.
  \item Click the row to open it and read the list of missing garments.
  \item \textbf{Physically look for them first.} Check the sending section, the
        floor, and the next trolley.
  \item If they truly cannot be found, press \btn{Close variance}.
  \item Type what happened. The system will not accept fewer than five characters
        --- write a real explanation, not a full stop.
  \item Press \btn{Close variance} to confirm.
\end{steps}

\shotsmall{75-transfers-doc}{A transfer document. Everything sent, everything
received, and exactly what is missing.}{fig:sup-doc}

\begin{plainwords}
Closing a variance does not delete the garments. They are put \textbf{on hold} in
the department that sent them, so they stop counting as ``in transit'' and stop
distorting everybody's figures. If one turns up next week, it is still in the
system with its whole history.
\end{plainwords}

\begin{stopbox}
Your name and the reason you typed are stored permanently against that document
and are visible in the audit trail. Closing a variance is an accounting decision,
so treat it as one.
\end{stopbox}

\section{Cancelling a handover sent by mistake}

If a batch was dispatched in error and the receiving department has \emph{not}
started receiving it:

\begin{steps}
  \item Open \menu{Transfers} for your section.
  \item Find the document under \emph{Dispatched, not yet received}.
  \item Press \btn{Cancel} and give a reason.
\end{steps}

All the garments return to your section immediately. Once even one piece has been
received, cancelling is no longer possible --- receive the rest and send it back
as a normal transfer instead.

\section{Watching the whole flow}

Open \menu{Dashboard}. Even if you only run one department, the flow board tells
you whether your section is about to be flooded or starved.

\shot{10-admin-dashboard}{The dashboard. Every department, what it is holding, and
what needs attention.}{fig:sup-dash}

Look at three things in particular:

\begin{bullets}
  \item \textbf{The card for the department before yours.} A large number there
        means work is coming.
  \item \textbf{``Inbound''} on your own card --- batches dispatched to you that
        nobody has received yet.
  \item \textbf{Transfers needing attention} --- counts that did not tally
        anywhere in the plant.
\end{bullets}

The coloured bar at the bottom of each card is the ageing profile: green is fresh
work, red has been sitting more than three days.

\section{Building a report}

Reports are built the same way by every role that has the screen, so the
walkthrough lives in its own chapter: see Chapter~\ref{ch:reports}. As a
supervisor the datasets you will use most are \emph{Articles / WIP} for what your
section is holding and \emph{Transfers} for handovers that have not settled.

\section{Fixing a garment that is genuinely stuck}

Occasionally a garment ends up recorded in the wrong place --- usually because a
trolley was moved before it was scanned.

\begin{steps}
  \item Open \menu{Trace Article} and find the garment by tag or serial number.
  \item Press \btn{Correct section}.
  \item Choose where it really is, and type why the correction is needed.
\end{steps}

\shot{99-trace-article}{The full life story of one garment: every scan, in order,
with who did it.}{fig:sup-trace}

If a tag is physically damaged, \btn{Replace damaged tag} lets you fit a new one.
The garment keeps its identity and its entire history; only the tag changes.

\begin{warnbox}
Both actions require a written reason and are recorded against your name. They
are for genuine data problems, not for making an awkward number go away.
\end{warnbox}
